\subsection{赋值环}

定理 1. 设 K 是域，V 是 K 的子环。下列条件等价：

\begin{enumerate}
\def\labelenumi{(\alph{enumi})}
\item
  V 是 K 的局部子环所成之集（该集以 A 与 B 之间的关系“B 支配 A”排序）的极大元。
\item
  存在代数闭域 L 与从 V 到 L 的同态 h，它在从 K 的子环到 L 的同态所成之集（该集以 f 与 g 之间的关系“g 是 f 的扩张”排序）中是极大的。
\item
  若 \(x \in \mathbf{K} - \mathbf{V}\)，则 \(x^{-1} \in \mathbf{V}\)。
\item
  \(\mathbf{V}\) 的分式域是 \(\mathbf{K}\)，且 \(\mathbf{V}\) 的主理想所成之集按包含关系是全序的。
\item
  V 的分式域是 K，且 V 的理想所成之集按包含关系是全序的。
\end{enumerate}

我们通过证明下列蕴涵关系来证明本定理

\[
(a) \Rightarrow (b) \Rightarrow (c) \Rightarrow (d) \Rightarrow (e) \Rightarrow (a).
\]

设 (a) 成立。则 V 是局部环。设 L 是剩余域 \(\kappa(V)\) 的一个代数闭包，h 是从 V 到 L 的典范同态。设 \(V'\) 是 K 的包含 V 的子环，\(h'\) 是从 \(V'\) 到 L 的延拓 h 的同态。若 \(p'\) 是 \(h'\) 的核，则 \(p' \cap V = m(V)\)；于是（第 1 节例 2）\(V_{p'}'\) 支配 V，这蕴涵 \(V_{p'}' = V\) 及 \(V' = V\)。故 (b) 成立。

设 (b) 成立。设 L 是代数闭域，h 是从 V 到 L 的同态；设 h 在从 K 的子环到 L 的同态所成之集中是极大的；设 p 是 h 的核。由于 \(h(V - p)\) 的元素都在 L 中可逆，h 可以延拓为从 \(V_{p}\) 到 L 的同态（第二章 §2 第 1 节命题 1）；于是 \(V = V_{p}\)，这表明 V 是局部环且 p 是它的极大理想。设 x 是 K 的非零元素；须证元素 x、\(x^{-1}\) 中至少有一个属于 V，也就是（依据 h 的极大性）h 可以延拓到 V{[}x{]} 或 \(V[x^{-1}]\)。若 x 在 V 上整，这由第五章 §2 第 1 节定理 1 的推论 4 得出。若 x 不在 V 上整，我们将使用下面的引理：

引理 1. 设 A 是环 B 的子环，x 是 B 的不在 A 上整的元素；则分式环 \(B_{x}\)（第二章 § 5 第 1 节）不约化为 0，且在子环 A{[}1/x{]}（\(B_{x}\) 的）中存在包含 1/x 的极大理想；此外，若 M 是这些极大理想中的任何一个，则 M 在 A 中的原像是极大理想。

由于 \(x\) 不在 \(A\) 上整，它不是幂零的，因此 \(B_x \neq 0\)；此外，\(x / 1 \notin A[1 / x]\)，否则将有一个形如

\[
x / 1 = a _ {0} / 1 + a _ {1} / x + \dots + a _ {n} / x ^ {n}
\]

的关系（其中 \(n \geqslant 0\)，且 \(a_{i} \in \mathbf{A}\)（对 \(0 \leqslant i \leqslant n\)）），它等价于

\[
x ^ {n + h} - a _ {0} x ^ {n + h - 1} - a _ {1} x ^ {n + h - 2} - \dots - a _ {n} x ^ {h - 1} = 0
\]

（其中 \(h \geqslant 1\) 取适当的值）；但这样一个关系将蕴涵 \(x\) 在 \(A\) 上整，与假设相反。于是，\(A[1/x]\) 的包含 \(1/x\) 的极大理想的存在性由 \(1/x\) 在 \(A[1/x]\) 中不可逆这一事实得出（《代数》第一章 § 8 第 7 节定理 2）。

于是设 \(\mathfrak{M}\) 是 A{[}1/x{]} 的包含 1/x 的一个极大理想；设 \(\phi: \mathrm{A} \to \mathrm{A}[1/x]\)，\(p: \mathrm{A}[1/x] \to \mathrm{A}[1/x]/\mathfrak{M}\) 是典范同态；则

\[
p (\mathrm{A} [ 1 / x ]) = p (\phi (\mathrm{A})) [ p (1 / x) ] = p (\phi (\mathrm{A}))
\]

因为 \(p(1/x) = 0\)；这证明 \(p(\phi(A))\) 是域，从而原像 \(\overline{\phi}^{-1}(M)\) 是 A 的极大理想。

我们以 A = V、B = K 应用此引理；于是存在 \(V[x^{-1}]\) 的包含 \(x^{-1}\) 的极大理想 M，且 \(M \cap V\) 是 V 的极大理想；由于 V 是局部的，有 \(M \cap V = p\)；用 f 表示从 \(V[x^{-1}]\) 到 \(V[x^{-1}] / M\) 上的典范同态，则 \(f(x^{-1}) = 0\)，从而 \(V / p = f(V) = f(V[x^{-1}])\)；由于 h 通过取商定义了从 V/p 到 L 的单同态 \(\bar{h}\)，\(\bar{h} \circ f\) 是从 \(V[x^{-1}]\) 到 L 的延拓 h 的同态。故 (c) 成立。

现在设 (c) 成立。显然 K 是 V 的分式域。设 a、b 是 V 的元素，且 \(Va \notin Vb\)；我们证明 \(Vb \subset Va\)。当 b = 0 时这是真的；否则关系 \(a \notin Vb\) 蕴涵 \(b^{-1}a \notin V\)，从而由 (c) 有 \(a^{-1}b \in V\)，因此 \(Vb \subset Va\)。故 (d) 成立。

设 (d) 成立。设 \(\mathfrak{a}\)、\(\mathfrak{b}\) 是 \(\mathbf{V}\) 的理想，且 \(\mathfrak{a} \notin \mathfrak{b}\)。存在 \(a \in \mathfrak{a}\) 使 \(a \notin \mathfrak{b}\)。对所有 \(b \in \mathfrak{b}\)，有 \(a \notin \mathbf{V}b\)，从而 \(Va \notin Vb\)，因此 \(Vb \subset Va \subset \mathfrak{a}\)（由 (c)），即 \(b \in \mathfrak{a}\)。于是 \(\mathfrak{b} \subset \mathfrak{a}\)，这表明条件 (e) 满足。

最后设 (e) 成立。由于 V 有极大理想，它只有一个，从而是局部环。设 \(V'\) 是 K 的支配 V 的局部子环，x 是 \(V'\) 的非零元素；记 \(x = ab^{-1}\)，其中 \(a \in V\)、\(b \in V\)。理想 Va、Vb 之一含于另一个。若 Va ⊂ Vb，则 \(x \in V\)。若 Vb ⊂ Va，则 \(x^{-1} \in V\)；由于理想 \(V'm(V)\) 不包含 1（第 1 节命题 1），\(x^{-1} \notin m(V)\)，于是仍得 \(x \in V\)，因为 \(V\) 是局部的。因此 \(V'\) 的每个元素都属于 \(V\)；我们断定 (a) 成立。

定义 2. 在定理 1 的记号下，若等价条件 (a)、(b)、(c)、(d)、(e) 成立，则称 V 是域 K 的赋值环。若一个环是整环且是它的分式域的赋值环，则称之为赋值环。

定理 2. 设 K 是域，h 是从 K 的子环 A 到代数闭域 L 的同态。则存在 K 的赋值环 V 与一个同态

\(h^{\prime}\) 从 \(\mathbf{V}\) 到 \(\mathbf{L}\)，使得 \(\mathbf{V}\) 包含 \(\mathbf{A}\)，\(h^{\prime}\) 延拓 \(h\)，且 \(h^{\prime}(0) = \mathfrak{m}(\mathbf{V})\)。

设 \(\mathfrak{S}\) 是从 K 的子环到 L 的同态所成之集，以扩张关系排序。这个集是归纳的；因为若 \((h_{\alpha})_{\alpha \in I}\) 是 \(\mathfrak{S}\) 的元素的一个非空全序族，\(B_{\alpha}\) 是 \(h_{\alpha}\) 的定义环，则诸 \(B_{\alpha}\) 构成 K 的子环的一个全序族，它们的并 B 因而是 K 的子环；于是存在唯一的从 B 到 L 的映射 \(\bar{h}\) 延拓诸 \(h_{\alpha}\)（《集合论》第二章 § 4 第 6 节命题 7），并且立即可见 \(\bar{h}\) 是从 B 到 L 的同态。Zorn 引理于是表明存在一个极大元 \(h'\)，它属于 \(\mathfrak{S}\) 且延拓 \(h\)。\(h'\) 的定义环 V 是 K 的赋值环（定理 1）；若 \(p\) 是 \(h'\) 的核，则 \(h'\) 可以延拓为从 \(V_p\) 到 L 的同态（第二章 § 2 第 1 节命题 1），于是 \(V_p = V\) 且 \(p = m(V)\)。

推论. 域 K 的每个局部子环 A 至少被 K 的一个赋值环所支配。

把定理 2 应用于典范同态 \(h\)，它从 \(\mathbf{A}\) 映到一个代数闭包 \(\mathbf{L}\)（\(\mathbf{A} / \mathfrak{m}(\mathbf{A})\) 的）中。

